% Future work section for OrbitGraph paper.

\section{Future Work}
\label{sec:future}

Our study demonstrates that graph-centric SDN can eliminate GSL handover black
holes on emulated constellations up to 102 nodes, but several limitations
suggest extensions in experimental methodology, scenario coverage, and
production-scale deployment.

The largest gap between OrbitGraph's algorithmic cost and its observed
control-plane time is kernel route install: pushing routes through per-container
remote execution dominates handover and damage events
(Section~\ref{sec:analysis-control}).
Replacing this harness with a production-style southbound, such as gRPC or
NETCONF to on-board agents with compressed delta messages, would isolate the
true cost of centralized programming from container overhead and likely improve
damage recovery, where install latency currently outweighs global recomputation.
Our synchronous StarryNet loop blocks during route install, inflating outage
measurements for reactive events. Running the controller asynchronously and
timing probes against control-plane timestamps would give a fairer comparison
with distributed protocols, whose reconvergence also overlaps forwarding.
The present campaign uses ICMP probes between two ground stations. Adding
sustained TCP or QUIC flows, multipath traffic matrices, and application-level
quality-of-service metrics would test whether zero handover outage translates
into throughput and goodput gains under realistic load, not only reachability.

Each run already exposes several handover events, which we classify by
ground-station coverage to separate routing behavior from physical connectivity.
Mega-constellations experience continuous contact churn. Emulating multi-hour
windows with dozens of handovers per ground station would reveal whether
incremental baselines stay consistent and whether error accumulation in address
caches or route-key bookkeeping appears over time. Denser constellations would
also close the coverage gaps that a sparse grid exposes at a fixed ground-station
latitude, isolating control-plane behavior from satellite visibility.
Our single autonomous system uses two fixed ground stations and LeastDelay
routing. Extending to multi-AS topologies with BGP peering, or to ground-relay
architectures that offload long-haul traffic to terrestrial
fiber~\cite{handley:ground-relays:2019}, would test OrbitGraph where path
selection crosses administrative boundaries and edge weights are not delay alone.
OrbitGraph proactive handover is complementary to distributed protocols that
exploit ephemeris. A head-to-head study against OPSPF~\cite{pam:opspf:2019} or
segment routing on the same topologies would clarify when centralized install
wins over orbit-predictive link-state flooding, and whether hybrid designs, an
OSPF data plane with SDN-triggered make-before-break, capture most of the
benefit with less operational coupling.
Future scenarios could include capacity-aware weighting beyond delay,
time-varying demand hotspots, partial link degradation such as elevated loss
without total failure, and targeted ISL failures rather than whole-satellite
blackouts.
Section~\ref{sec:analysis-damage} shows OrbitGraph lagging OSPF after
30\% satellite loss because updates are reactive and bulk. Promising directions
include precomputed backup next-hops for known failure domains, staged install
that restores connectivity to ground before refining intra-orbit paths, and soft
damage detection that triggers recomputation before total link loss.

Container emulation allocates one Linux network stack per node. Our workstation
campaign tops out near 102 nodes before memory and CPU become limiting factors,
consistent with observations in the StarryNet literature~\cite{lai2023starrynet}.
Reaching Starlink-class scale (thousands of satellites) requires a tiered
strategy: partition the constellation into orbital planes or regions and run
OrbitGraph per partition with border-route summarization at domain
gateways~\cite{gregory:satellite-routing:2022}; shard containers across hosts
using StarryNet's distributed deployment mode; or combine packet-level emulation
on a representative subgraph with flow-level simulators such as
Hypatia~\cite{kassing2020exploring} or StarPerf~\cite{lai2020starperf} for
full-constellation stress tests, using the same $G_t$ abstraction in both tiers.
Our grids use Walker-style $O{\times}S$ layouts with Starlink-inspired altitude
and inclination but not live two-line element feeds.

Closing the install-time gap, broadening workloads and failure models, and
climbing the scaling ladder through hierarchy and hybrid simulation are the
main levers to move from proof-of-concept to operator-relevant deployment.
The core idea, treating the LEO network as a sequence of graphs, computing
shortest paths centrally, and timing delta installs against known topology
events, remains unchanged; future work fills in the engineering surround
that our current emulation study deliberately simplified.
